\section{总结与延伸}

\subsection{全课知识图谱}

本课建立了从神经网络到反向传播的完整认知链：

\begin{center}
\begin{tikzpicture}[node distance=0.8cm, every node/.style={draw, rounded corners, minimum width=3.5cm, minimum height=0.7cm, font=\small}]
\node (linear) {线性分类器};
\node (nn) [below=of linear] {神经网络};
\node (act) [below=of nn] {激活函数};
\node (cg) [below=of act] {计算图};
\node (bp) [below=of cg] {反向传播};
\node (ad) [below=of bp] {自动微分框架};
\draw[->, thick] (linear) -- (nn);
\draw[->, thick] (nn) -- (act);
\draw[->, thick] (act) -- (cg);
\draw[->, thick] (cg) -- (bp);
\draw[->, thick] (bp) -- (ad);
\node[draw=none, right=1cm of linear, text width=5cm, font=\scriptsize] {$f = W\mathbf{x}$，线性不可分};
\node[draw=none, right=1cm of nn, text width=5cm, font=\scriptsize] {$f = W_2 \sigma(W_1 \mathbf{x})$，非线性决策边界};
\node[draw=none, right=1cm of act, text width=5cm, font=\scriptsize] {ReLU, GELU, SiLU, Sigmoid};
\node[draw=none, right=1cm of cg, text width=5cm, font=\scriptsize] {前向传播，每个操作是一个节点};
\node[draw=none, right=1cm of bp, text width=5cm, font=\scriptsize] {链式法则，上游$\times$局部梯度};
\node[draw=none, right=1cm of ad, text width=5cm, font=\scriptsize] {PyTorch, TensorFlow};
\end{tikzpicture}
\end{center}

\subsection{关键 Takeaways}

\begin{importantbox}{六条核心原则}
\begin{enumerate}
    \item \textbf{非线性是神经网络的灵魂}：没有激活函数，多层网络等价于单层线性模型
    \item \textbf{反向传播 = 链式法则的系统化应用}：将复杂函数分解为简单操作，逐个求导
    \item \textbf{每个节点只需关心局部}：计算局部梯度，乘以上游梯度，传给下游
    \item \textbf{加法分发、乘法交换、max 路由}：三种基本操作的梯度模式是所有复杂网络的构建块
    \item \textbf{维度匹配是最好的调试工具}：梯度矩阵的形状必须与对应参数矩阵相同
    \item \textbf{自动微分解放了研究者}：理解反向传播原理后，实践中由框架自动完成
\end{enumerate}
\end{importantbox}

\subsection{拓展阅读}

\begin{itemize}
    \item Rumelhart, Hinton \& Williams, Learning Representations by Back-propagating Errors (1986): \url{https://www.nature.com/articles/323533a0} —— 反向传播的经典论文
    \item Goodfellow, Bengio \& Courville, Deep Learning, Chapter 6: \url{https://www.deeplearningbook.org/} —— 深度前馈网络详解
    \item Stanford CS231N 作业 1: \url{http://cs231n.stanford.edu/} —— 手动实现反向传播的练习
    \item PyTorch Autograd 文档: \url{https://pytorch.org/docs/stable/autograd.html} —— 自动微分机制详解
    \item Andrej Karpathy, micrograd: \url{https://github.com/karpathy/micrograd} —— 一个极简自动微分引擎，适合理解反向传播
\end{itemize}

\end{document}
